% Actualizacion 2024: Logotipos de la UPV


\documentclass[12pt]{article}
\usepackage[left=2.5cm,top=2.0cm,right=2.5cm,bottom=3.0cm]{geometry}
\usepackage[utf8]{inputenc}
\usepackage[spanish]{babel}
%\usepackage[english, american]{babel}
%\usepackage[spanish,es-tabla]{babel}
\usepackage[linguistics]{forest}
\usepackage{amssymb, amsmath, amsbsy} % simbolitos
\usepackage{longtable} % para tablas largas
\usepackage{graphicx}
\usepackage{fancyhdr}
\usepackage{xcolor}
\usepackage{multirow}
\usepackage{listings}
\usepackage{caption}
\usepackage{subcaption}
%\usepackage{parskip}
\usepackage[skip=12pt plus1pt]{parskip}
\usepackage{pdfpages} % Incluir PDF en documento en LATEX
\usepackage{verbatim} % comentarios
\usepackage{algpseudocode}
\usepackage{algorithm}
\usepackage{pdflscape}
\usepackage{multirow}
\usepackage{afterpage}
\usepackage{array,booktabs,ragged2e}
%\newcolumntype{R}[1]{>{\RaggedLeft\arraybackslash}p{#1}}
\newcolumntype{L}[1]{>{\raggedright\let\newline\\\arraybackslash\hspace{0pt}}m{#1}}
\newcolumntype{C}[1]{>{\centering\let\newline\\\arraybackslash\hspace{0pt}}m{#1}}
\newcolumntype{R}[1]{>{\raggedleft\let\newline\\\arraybackslash\hspace{0pt}}m{#1}}

\floatname{algorithm}{Algoritmo}
\renewcommand{\listalgorithmname}{Lista de algoritmos}
\renewcommand{\algorithmicrequire}{\textbf{Entrada:}}
\renewcommand{\algorithmicensure}{\textbf{Salida:}}

% Comentario con respecto a las referencias:
% Por defecto este documento utiliza el formato IEEEtr para formatear las referencias. Si algun asesor requiere formato APA, solo comente la siguiente linea y descomente la linea debajo. Tambien para que las citas funcionen de manera adecuada, el paquete label debe tener las opciones mostradas, o en su defecto, el año en el listado no lo muestra correctamente \usepackage[english, american]{babel}

% Para utilizar el formato de citas IEEE y comentar los dos parrafos siguientes
\usepackage[backend=bibtex,sorting=none]{biblatex}
% Para utilizar el formato APA, sugiero comentar la linea anterior y descomentar las dos proximas lineas
%\usepackage[backend=biber,style=apa]{biblatex}
%\DeclareLanguageMapping{english}{american-apa}

\makeatletter
\DefineBibliographyExtras{spanish}{%
  \setcounter{smartand}{1}%
  \let\lbx@finalnamedelim=\lbx@es@smartand
  \let\lbx@finallistdelim=\lbx@es@smartand
}
\renewbibmacro*{name:delim:apa:family-given}[1]{%
  \ifnumgreater{\value{listcount}}{\value{liststart}}
    {\ifboolexpr{
       test {\ifnumless{\value{listcount}}{\value{liststop}}}
       or
       test \ifmorenames
     }
       {\printdelim{multinamedelim}}
       {\lbx@finalnamedelim{#1}}}
    {}}
\makeatother


% Estas lineas permiten romper los hipervinculos muy largos en las referencias!!!!
\setcounter{biburllcpenalty}{7000}
\setcounter{biburlucpenalty}{8000}
\addbibresource{x.bib} % ARCHIVO DE BIBLIOGRAFÍA


%\usepackage{url}
\usepackage[bookmarks=true,breaklinks=true,bookmarksopen=false,colorlinks=true,linkcolor=blue]{hyperref}
\usepackage[hyphenbreaks]{breakurl}
% Regla que define explicitamente que caracteres rompen los hipervinculos para separar las lineas
%https://es.overleaf.com/11089898rhgykrqyqytx
% Actualiza en automático la fecha de las citas de internet a la fecha de la compilación del documento
\usepackage{datetime}
\newdateformat{specialdate}{\twodigit{\THEDAY}-\twodigit{\THEMONTH}-\THEYEAR}
%\newdateformat{specialdate}{\twodigit{\THEDAY}-\THEYEAR}
\date{\specialdate\today}

\newcommand{\HRule}{\rule{\linewidth}{0.25mm}}


% CONSTANTES NECESARIAS PARA EL DOCUMENTO ---> MODIFIQUEN A SU CRITERIO
\newcommand{\ncarrera}            {Ingeniería en Tecnologías de la Información}
\newcommand{\nasesorinstitucional}{Dr. Said Polanco Martagón}



% ALUMNOS
\newcommand{\NombreAlumno}{Erika Daniela Mallozzi Martínez y Heriberto Geovanny Nava Lopez}
%Hombres cambien LA por EL
\newcommand{\elolaNombreAlumno}{la}  
\newcommand{\OA}           {a}  %Hombres cambien A por O
\newcommand{\Matricula}{2130128 y 2130213}





\newcommand{\NombreProyecto}{Automatización de Clasificación de Limones por Color.}
\newcommand{\fechacarta}{26 de Abril de 2021}
\newcommand{\ncuatrimestre}{Mayo-agosto 2021}
\newcommand{\nevalador}{Dr. Said Polanco Martagón}
\newcommand{\FechaExposicion}{11 de Agosto de 2021}
\newcommand{\HoraExposionFormatoVenticuatroHoras}{10:00}
\newcommand{\elolaNombreEmpresa}{la}  %Si la empresa es femenino (por ejemplo universidad, usen la) o masculino (el instituto) pongan el
\newcommand{\organismoreceptor}   {Universidad Politécnica de Victoria}

% NOTA: Dos diagonales juntas (\\) indican un saldo de linea. En este caso particular hay 2 (el titulo se ajusta a tres lineas, porque es muy largo. Hacer las adecuaciones pertinentes
\newcommand{\NombreProyectoheader}     {Automatización de Clasificación de Limones por Color}
\newcommand{\nasesorempresaria}   {Dr. Said Polanco Martagón}
\newcommand{\fechaPortada}               {Octubre de 2024}





\newcommand{\separacionCorta}{0.0cm}
\newcommand{\separacionLarga}{0.5cm}

\usepackage[overload]{textcase}
\newcommand{\iemph}[1]{\MakeTextUppercase{#1}}

\pagestyle{fancy}
\headheight 45pt
\fancyhead{} % Clear all header fields
\fancyhead[L]{\includegraphics[height=1.00cm]{UTYP.png}}%
\fancyhead[C]{\begin{center}\NombreProyectoheader\end{center}}%
%\fancyhead[R]{\includegraphics[height=1.5cm]{LogoUPV_2019.png}}%
\fancyhead[R]{\includegraphics[height=1.25cm]{LogoUPV_2023.png}}%
\fancyfoot[R]{\thepage} % Clear all footer fields 
\fancyfoot[C]{}
\fancyfoot[L]{}

\DefineBibliographyStrings{english}{%
  references = {Referencias},% replace "references" with "bibliography"  for `book`/`report`
}

\addto\captionsenglish{%
  \renewcommand{\figurename}{Figura}%
  \renewcommand{\tablename}{Tabla}%
} 

\usepackage{wallpaper}
 
 
%\renewcommand{\figurename}{Figura}
%\renewcommand{\tablename}{Tabla}

 
\begin{document}

%-----------------------------------------------------------------------------------------------------------------
% PAGINA 1 - PORTADA
\setcounter{page}{1}
\pagenumbering{roman}
\thispagestyle{empty}

\begin{center}

\begin{tabular}{cp{5cm}c}
\includegraphics[height=2.25cm]{UTYP.png} & 
& \includegraphics[height=2.25cm]{LogoUPV_2023.png}   \\
\end{tabular}

\Large \textbf{UNIVERSIDAD POLITÉCNICA DE VICTORIA}
\vspace{0.5cm}
\hrule
\vspace{0.1cm} 
\hrule
\vspace{0.5cm}


%\HRule \\[\separacionCorta]
\textbf{\iemph{\NombreProyecto}} \\[\separacionLarga]
%\Large \textbf{TESINA}
%\HRule \\[\separacionLarga]
T E S I N A \\

\textbf{\iemph{\ncarrera}} \\[\separacionLarga]

PRESENTA: \\[\separacionCorta]
%\textbf{\Capitalize{\NombreAlumno}\\[\separacionLarga]
\textbf{\iemph{\NombreAlumno}}\\[\separacionLarga]
%EN CUMPLIMIENTO DE \\[\separacionCorta]
%LA ESTADÍA DE LA CARRERA DE \\[\separacionCorta]

\vspace{3cm}

DIRECTOR \\[\separacionCorta]
\textbf{\iemph{\nasesorinstitucional}} \\[\separacionCorta]



ORGANISMO RECEPTOR \\[\separacionCorta]
\textbf{\iemph{\organismoreceptor}} \\[\separacionLarga]

\end{center}
\begin{flushright}
\iemph{Ciudad Victoria, Tamaulipas, \fechaPortada}
\end{flushright}

\HRule 




%-----------------------------------------------------------------------------------------------------------------

%-----------------------------------------------------------------------------------------------------------------
% PAGINA 7 - INDICE

\clearpage
\addcontentsline{toc}{section}{Índice}
\renewcommand\contentsname{Índice}
\tableofcontents

%-----------------------------------------------------------------------------------------------------------------
% CAPITULOS


\clearpage
\pagenumbering{arabic}
\setcounter{page}{1}
%\input{Capitulo1.tex}

\clearpage
%\input{Capitulo2.tex}

\clearpage
\section{Introducción}

Cuando hablamos de automatización, es fundamental considerar las necesidades que se pretenden satisfacer en cualquier industria. Los empresarios, especialmente en el ámbito de la agricultura, buscan constantemente maneras de agilizar y modernizar sus procesos de producción. Esto responde a la creciente demanda de eficiencia, precisión y ahorro de recursos. En este contexto, es inevitable hacer referencia a la Revolución Industrial, un período que cambió radicalmente la sociedad, la economía y la tecnología. Esta revolución dio origen a la creación de máquinas capaces de realizar tareas que antes dependían exclusivamente de la mano de obra humana, y sentó las bases para los avances tecnológicos que hoy conocemos.

La automatización en la agricultura no es una excepción. De hecho, en la actualidad se concentra en la creación de máquinas inteligentes diseñadas para optimizar una amplia gama de actividades, desde el cultivo y riego hasta la recolección y clasificación de productos. Entre estas actividades, la clasificación de alimentos ocupa un lugar clave. Este proceso es esencial para asegurar que los productos que llegan al consumidor final cumplan con los estándares de calidad que el mercado exige, especialmente en un mundo donde la demanda de alimentos frescos y de alta calidad es cada vez mayor.

La importancia de la clasificación de alimentos está profundamente arraigada en nuestra cultura. Desde temprana edad, aprendemos la relevancia de una alimentación balanceada y de calidad, algo que se nos inculca a través de guías como el "Plato del Buen Comer". A lo largo de la vida, el acto de seleccionar alimentos en mercados o tiendas se convierte en una práctica cotidiana que, aunque sencilla en apariencia, es crucial para garantizar que lo que consumimos sea saludable y nutritivo. Al elegir frutas y verduras, por ejemplo, no solo tomamos en cuenta su apariencia, sino también su frescura, madurez y estado general. Este proceso manual de selección, que en muchas familias se convierte en una tradición, es una forma intuitiva de asegurar la calidad de los alimentos, pero también es ineficiente y propenso a errores cuando se trata de grandes volúmenes en entornos industriales.

Un claro ejemplo de esta problemática se presenta en la producción de cítricos, específicamente en la clasificación de limones. La calidad de estos productos varía considerablemente dependiendo de la temporada y del tipo de limón. El limón verde, que tiene una temporada de septiembre a junio, se distingue por su acidez, ausencia de semillas y apariencia externa, características que lo diferencian de otros tipos de limones. Si bien es relativamente sencillo para un agricultor o seleccionador experimentado identificar estas diferencias visuales, este proceso manual es lento y está sujeto a errores humanos. La clasificación de grandes volúmenes de limones basándose únicamente en observaciones visuales no solo consume tiempo, sino que también puede resultar en productos clasificados incorrectamente, lo que afecta negativamente tanto la calidad del producto como la rentabilidad de los productores.

Es en este punto donde la automatización puede marcar una diferencia sustancial. La implementación de sistemas de clasificación automatizados, basados en parámetros como el tamaño, el color y otras características visuales, permite optimizar el proceso y reducir significativamente los tiempos de inspección. A través de la automatización, se logra una clasificación más precisa y consistente, lo que a su vez mejora la calidad de los productos que llegan al consumidor final. Los sistemas automatizados no solo seleccionan con mayor precisión los limones que cumplen con los estándares de calidad, sino que también eliminan aquellos que no son aptos para la venta, lo que minimiza las pérdidas y optimiza el uso de los recursos disponibles.

El uso de tecnologías avanzadas, como la visión artificial, es una de las soluciones más prometedoras en este campo. Esta tecnología permite a las máquinas "ver" y analizar las características visuales de los productos, clasificándolos con una precisión que supera a la del ojo humano. La visión artificial, combinada con algoritmos de aprendizaje automático, puede identificar de manera eficiente atributos como el color, el tamaño, la madurez y otros factores clave que determinan la calidad de un limón. Estas innovaciones no solo mejoran la precisión en la clasificación, sino que también aumentan la velocidad del proceso, lo que es fundamental en un entorno de producción agrícola a gran escala.

En este proyecto, el objetivo principal es desarrollar un sistema automatizado que utilice la visión artificial para clasificar limones según su calidad, basado en parámetros predefinidos de color y tamaño. Este sistema no solo promete aumentar la eficiencia del proceso de clasificación, sino que también busca reducir el margen de error humano y garantizar que solo los limones en perfecto estado lleguen al mercado. De esta manera, se pretende contribuir tanto a la sostenibilidad del sector agrícola como a la satisfacción del consumidor final, quien espera productos frescos y de alta calidad.

La implementación de este tipo de sistemas puede tener un impacto significativo en la industria agrícola. A medida que los productores adoptan estas tecnologías, no solo se optimizan los procesos internos, sino que también se mejora la competitividad en el mercado, al poder ofrecer productos de mejor calidad de manera más eficiente. La automatización, por tanto, no es simplemente una tendencia, sino una necesidad creciente para los productores que buscan mantenerse al día con las demandas del mercado y aprovechar al máximo los recursos disponibles.

\clearpage
\section{Antecedentes Teóricos}

La clasificación de productos agrícolas, como los cítricos, es crucial para garantizar la calidad del producto y su valor en el mercado. Con el aumento de la demanda de productos frescos y de alta calidad, la automatización de este proceso se ha convertido en una necesidad apremiante. Las tecnologías de visión artificial han emergido como herramientas clave en este ámbito, con investigaciones recientes que ofrecen perspectivas valiosas, pero que también presentan limitaciones.

\subsection{Importancia de la Clasificación en la Industria Agrícola}

La clasificación efectiva de los limones asegura su calidad y determina su valor comercial. Según Vargas et al. (2020), la adopción de sistemas automatizados puede mejorar la calidad del producto al reducir la variabilidad y los errores de clasificación asociados con los métodos manuales. Sin embargo, estas soluciones automatizadas aún pueden ser ineficaces si no utilizan técnicas avanzadas para la identificación de características específicas de los limones.

\subsection{Desarrollo de Tecnologías de Visión Artificial}

El avance en las tecnologías de visión artificial ha permitido la creación de sistemas de clasificación más eficientes. Guevara et al. (2020) destacan que los sistemas que utilizan algoritmos de aprendizaje profundo pueden procesar imágenes en tiempo real, mejorando la identificación de los limones según sus características visuales. No obstante, muchos de estos sistemas se centran en parámetros limitados, lo que puede afectar la precisión de la clasificación. Por ejemplo, algunos métodos no consideran adecuadamente las variaciones de color que indican el estado de madurez del fruto, lo que puede llevar a clasificaciones incorrectas y pérdidas de productos.

\subsection{Impacto en la Productividad y Sostenibilidad}

Las investigaciones también indican que la inversión en tecnologías automatizadas, aunque beneficiosa a largo plazo, no siempre es accesible para todos los productores. Vargas et al. (2020) señalan que, si bien la automatización mejora la productividad, muchos sistemas requieren una intervención constante de los operarios, lo que limita su eficacia. En contraste, el sistema propuesto en este estudio aborda directamente estas limitaciones al desarrollar una solución automatizada de clasificación de limones basada en visión artificial y colorimetría, utilizando un ESP32-CAM y un sensor de color.

\subsection{Solución Propuesta: Sistema Automatizado de Clasificación}

La propuesta consiste en un sistema automatizado de clasificación de limones que integra tecnologías avanzadas para evaluar la calidad de la fruta en función del color de la piel externa. A diferencia de los sistemas existentes que pueden ser limitados en su enfoque, este sistema se propone:

- Reducción significativa de ineficiencia: Se anticipa que el nuevo sistema reducirá el tiempo y el esfuerzo manual actualmente requeridos en la clasificación, minimizando la dependencia de operarios y permitiendo una clasificación más rápida y precisa.
  
- Mejor identificación de limones en mal estado: Gracias al uso de parámetros de color establecidos y a un algoritmo de procesamiento de imágenes que emplea técnicas de análisis de colorimetría y machine learning, se espera que el sistema logre identificar de manera precisa los limones en mal estado antes de su empaque. Esto no solo disminuirá la pérdida de productos, sino que también mejorará la calidad del producto final.

- Aumento en la satisfacción del consumidor: Al asegurar que únicamente los limones en su punto óptimo de madurez lleguen al mercado, se incrementará la satisfacción del consumidor final, lo que es fundamental para la competitividad en el sector agrícola.

\subsection{Objetivos del Sistema Automatizado}

Los objetivos específicos del sistema propuesto son:

1. Desarrollar un sistema automatizado que permita a los productores clasificar los limones según su calidad, reduciendo el error humano y optimizando el proceso de clasificación.
  
2. Diseñar algoritmos y módulos de procesamiento que clasifiquen los limones por color y calidad, mejorando la precisión en la evaluación visual.

3. Definir una metodología de evaluación de parámetros de color, que identifique con precisión la madurez y calidad de los limones mediante métricas estandarizadas.

4. Validar la fiabilidad del sistema mediante pruebas controladas y comparación de resultados con estándares de color existentes para limones.

5. Realizar pruebas de integración entre los componentes del sistema en un entorno de producción real, asegurando tiempos de respuesta rápidos y alta precisión en la clasificación de grandes volúmenes de limones.

Con este enfoque, se busca superar las limitaciones de las investigaciones anteriores, ofreciendo una solución más eficiente y adaptable a las necesidades del sector agrícola.


\clearpage
\section{Marco Teórico}

La clasificación de frutas en la industria agrícola es un proceso crucial para garantizar la calidad del producto, especialmente en mercados donde se priorizan estándares específicos de madurez y apariencia. En este proyecto, se aborda la automatización de la clasificación de limones a través de técnicas de visión artificial, utilizando la ESP32-CAM como elemento clave. A continuación, se describen los conceptos teóricos y metodológicos que fundamentan el desarrollo de este sistema.

\subsection{Teoría del Color}

El color es uno de los principales indicadores visuales de la madurez en los limones. La teoría del color proporciona un marco para analizar y clasificar los diferentes tonos que se presentan durante el ciclo de maduración del fruto. En este proyecto, el modelo de color \textbf{HSV} (Hue, Saturation, Value) se utiliza debido a su capacidad para separar la información de matiz (color) de la intensidad, lo que facilita la segmentación del color del limón en distintos rangos de madurez. Esto permite clasificar los frutos de acuerdo a niveles como verde, amarillo o maduro.

\subsection{Visión Artificial}

La visión artificial es una disciplina que involucra la interpretación de imágenes digitales para obtener información relevante. En el contexto de este proyecto, la visión artificial permite identificar patrones visuales en los limones mediante el análisis de las imágenes capturadas por la \textbf{ESP32-CAM}. El procesamiento de imágenes involucra pasos como la segmentación, detección de características y clasificación de objetos en la imagen. Este enfoque mejora la eficiencia de la selección, eliminando el factor subjetivo del proceso manual.

\subsection{Análisis de Textura y Color}

El análisis de textura y color es fundamental en la clasificación de frutas. En el caso de los limones, la textura de la cáscara y el color externo pueden ofrecer información clave sobre su madurez y calidad. La combinación de estos dos parámetros permite implementar un sistema de clasificación más robusto, utilizando la cámara ESP32-CAM para capturar imágenes y realizar una discriminación basada en características visuales.

\subsection{Teoría de Clasificación}

La clasificación de los limones en este proyecto sigue los principios de la teoría de clasificación supervisada. El objetivo es agrupar los limones en diferentes clases basadas en su color, lo que se realiza entrenando un modelo de machine learning para reconocer patrones en las imágenes capturadas. El sistema aprende de un conjunto de datos previamente etiquetado, y luego utiliza esta información para clasificar nuevas muestras de limones de manera autónoma.

\subsection{Metodologías}

Este proyecto sigue una metodología de \textbf{aprendizaje supervisado}, en la que se entrena un modelo con un conjunto de datos previamente etiquetado. Las imágenes de limones, clasificadas por su madurez, son procesadas para extraer información del color y la textura. Posteriormente, el sistema utiliza estos datos para predecir la clase de madurez de limones nuevos. 

Adicionalmente, se emplea una \textbf{metodología de procesamiento de imágenes digitales}, que involucra la transformación de las imágenes capturadas por la ESP32-CAM. Este proceso incluye la conversión de espacios de color (por ejemplo, de RGB a HSV), la segmentación de las áreas relevantes del limón y la eliminación de ruido, lo que permite una interpretación más precisa de las características visuales.

El desarrollo del sistema también se enmarca en una \textbf{metodología iterativa}, donde el proyecto se divide en fases. Cada fase incluye pruebas parciales y ajustes del modelo, lo que permite la mejora continua del sistema de clasificación a medida que se obtiene nueva retroalimentación.

\subsection{Métodos}

En cuanto a los métodos utilizados, se hace uso de diversas técnicas de machine learning y visión artificial para implementar el sistema de clasificación de limones:

\begin{itemize}
    \item \textbf{K-Vecinos más Cercanos (KNN)}: Se emplea el algoritmo de \textit{K-Nearest Neighbors} para clasificar los limones en diferentes categorías basadas en su color. Este método compara el color de los limones nuevos con los colores de un conjunto de limones previamente clasificados.
    
    \item \textbf{Redes Neuronales Convolucionales (CNN)}: Para una clasificación más avanzada, se utiliza el método de \textit{Redes Neuronales Convolucionales}, que aprende patrones complejos en las imágenes de los limones. Este enfoque permite identificar de manera más precisa las diferencias sutiles en la textura y color que indican la madurez del fruto.
    
    \item \textbf{Segmentación de Color}: El método de segmentación se utiliza para aislar las áreas relevantes de la imagen del limón, basado en su color. Se aplica un umbral en el espacio de color \textbf{HSV} para distinguir diferentes niveles de madurez, separando, por ejemplo, las regiones verdes de las amarillas.
    
    \item \textbf{K-Means Clustering}: Para explorar la agrupación no supervisada, se utiliza el algoritmo \textit{K-Means}, que agrupa los limones según su similitud en características visuales, sin necesidad de etiquetas previas. Esto es útil en las fases iniciales de la investigación para definir patrones visuales de madurez.
\end{itemize}



\clearpage
\section{Estado del Arte}

La clasificación automática de frutas ha sido un tema de investigación activo durante las últimas décadas debido a la creciente necesidad de optimizar los procesos de selección en la industria agrícola. En particular, el uso de visión artificial ha permitido importantes avances en la automatización de la clasificación de frutas, mejorando tanto la eficiencia como la precisión en la evaluación de calidad.

\subsection{Avances en visión artificial para la clasificación de frutas}

Durante los últimos cinco años, la visión artificial ha experimentado avances significativos, impulsados por el desarrollo de cámaras más asequibles y la mejora de algoritmos de procesamiento de imágenes. Investigaciones recientes han utilizado técnicas de procesamiento digital de imágenes, aprendizaje automático y redes neuronales convolucionales (CNN) para analizar parámetros externos de las frutas, como color, textura y forma. Estos estudios han demostrado la efectividad de la visión artificial para clasificar una amplia variedad de productos agrícolas, incluyendo manzanas, tomates, y cítricos, entre otros.

Por ejemplo, estudios como los de Xie et al. (2021) demostraron el uso de redes neuronales profundas para la clasificación de cítricos en función de su madurez, utilizando el color como uno de los principales criterios. Otros trabajos, como el de Zhang et al. (2020), han desarrollado modelos de clasificación multietiqueta que integran tanto el color como la textura para obtener una evaluación más precisa de la calidad del fruto.

En el caso específico de los limones, la clasificación por color ha sido objeto de interés debido a su importancia en la determinación del grado de madurez. Varias investigaciones han señalado que el color externo es uno de los indicadores más fiables de la madurez en cítricos, y los sistemas de visión artificial basados en este criterio han demostrado ser altamente eficientes para este propósito. No obstante, la mayoría de estos estudios utilizan cámaras de alta resolución y sistemas de procesamiento avanzados, lo que incrementa el costo y la complejidad del sistema.

\subsection{Uso de sistemas embebidos para la clasificación de frutas}

En los últimos años, los dispositivos de bajo costo como la ESP32-CAM han comenzado a atraer la atención de los investigadores debido a su capacidad para implementar sistemas de visión artificial en entornos con recursos limitados. A diferencia de los sistemas más costosos, la ESP32-CAM ofrece una solución accesible para pequeños productores agrícolas que no pueden invertir en infraestructura tecnológica avanzada. Sin embargo, las investigaciones sobre su aplicación en la clasificación de frutas aún son limitadas.

A pesar de su bajo costo, estudios recientes han demostrado que la ESP32-CAM, combinada con algoritmos de procesamiento optimizados, puede ofrecer resultados satisfactorios en tareas de clasificación basadas en color. Un ejemplo es el trabajo de García et al. (2022), donde se utilizó la ESP32-CAM para clasificar manzanas por color con una precisión considerable, lo que sugiere su potencial para la clasificación de otros frutos como los limones.

\subsection{Desafíos y oportunidades}

A pesar de los avances en la clasificación automática de frutas mediante visión artificial, aún existen varios desafíos. La variabilidad de las condiciones de iluminación, la presencia de defectos físicos no relacionados con el color, y las limitaciones en la capacidad de procesamiento de los sistemas embebidos son factores que pueden afectar la precisión de los modelos de clasificación. Sin embargo, estos desafíos también presentan oportunidades para la innovación, especialmente en el contexto de sistemas accesibles como el que se plantea en este proyecto.

El presente trabajo se sitúa en un contexto en el que se busca democratizar el uso de tecnologías avanzadas en la industria agrícola, permitiendo a pequeños productores acceder a soluciones tecnológicas eficientes y de bajo costo. Al utilizar la ESP32-CAM como base, este proyecto pretende aportar una solución accesible y eficiente que, aunque centrada en la clasificación por color, podría expandirse para incluir otros parámetros de calidad en futuras investigaciones.

Aunque existen desarrollos recientes significativos en la automatización de la clasificación de frutas mediante visión artificial, este proyecto se diferencia al enfocarse en un dispositivo accesible como la ESP32-CAM. Se espera que, al optimizar su uso para la clasificación de limones por color, se contribuya a llenar el vacío existente en la investigación aplicada a pequeños productores con recursos limitados.


\clearpage
\section{Descripción del Problema}

La clasificación de limones sigue siendo, en gran medida, un proceso manual que enfrenta numerosas limitaciones. A pesar de los avances en la automatización en otras áreas de la agricultura, la clasificación de frutas como el limón todavía depende en gran parte de trabajadores que realizan evaluaciones visuales, lo que da lugar a varios problemas relacionados con la inconsistencia, la ineficiencia y los altos costos. Esta actividad manual conlleva errores humanos frecuentes que impactan directamente en la calidad del producto y, por ende, en la satisfacción de los clientes.

Un problema evidente en el proceso de clasificación manual es la falta de precisión y uniformidad. Los trabajadores encargados de clasificar los limones se basan en su experiencia y percepción personal para evaluar características como el tamaño, el color y la madurez de la fruta. Sin embargo, esta subjetividad genera inconsistencias en los resultados. Un trabajador puede considerar que un limón cumple con los estándares de calidad, mientras que otro puede clasificar el mismo limón de manera diferente. La falta de criterios estandarizados afecta negativamente la calidad general del producto y puede llevar a una mezcla de frutas de distintos grados de madurez y calidad en el mismo lote, lo que genera descontento entre los consumidores.

Otro desafío significativo es la incapacidad del método manual para manejar grandes volúmenes de producto de manera eficiente. A medida que la demanda de limones aumenta tanto a nivel local como internacional, los productores enfrentan la presión de clasificar mayores cantidades en menos tiempo. Sin embargo, el proceso manual es lento y requiere un esfuerzo considerable. Esto no solo retrasa la distribución de los productos, sino que también limita la capacidad de los productores para cumplir con los plazos establecidos. En un entorno competitivo, donde la frescura y calidad del producto son factores determinantes para el éxito, los productores que no puedan acelerar sus procesos de clasificación corren el riesgo de quedarse atrás.

La falta de estandarización en los criterios de calidad es otro aspecto preocupante. Cada trabajador puede aplicar diferentes parámetros al clasificar los limones, lo que resulta en una clasificación irregular. Esta variabilidad no solo afecta la calidad del producto, sino que también genera problemas en la cadena de suministro, ya que los distribuidores y minoristas confían en productos consistentes que cumplan con los estándares establecidos. Cuando los limones de menor calidad son mal etiquetados como productos de alta calidad, no solo se compromete la satisfacción del cliente, sino que también se producen pérdidas económicas tanto para los productores como para los distribuidores.

Uno de los mayores riesgos de la clasificación manual es la propensión a cometer errores, especialmente en situaciones de alta presión o cuando los trabajadores están cansados. Según Sánchez et al. (2022), hasta un 15\% de los limones clasificados manualmente son rechazados por defectos no detectados. Esta cifra representa una pérdida significativa, tanto en términos de ingresos como de calidad. Los errores humanos no solo afectan la reputación de los productores, sino que también generan una cantidad considerable de desperdicio. En un contexto donde la sostenibilidad y la reducción de residuos son prioritarias, estas pérdidas son inaceptables.

Los costos operativos asociados a la clasificación manual son igualmente alarmantes. La dependencia de mano de obra calificada conlleva altos costos salariales y de capacitación. Según estimaciones de López y Pérez (2023), los gastos de clasificación pueden representar hasta un 20\% del costo total de producción. Este margen de costos afecta directamente la rentabilidad de los productores, quienes deben encontrar formas de reducir gastos para mantener su competitividad en el mercado. En este sentido, la clasificación manual no es un método sostenible ni eficiente a largo plazo.

La creciente preocupación por la sostenibilidad también debe tenerse en cuenta. Ramírez y Torres (2020) señalan que hasta un 30\% de las frutas y verduras se desperdician debido a ineficiencias en la clasificación y el empaquetado. La incapacidad de detectar limones en mal estado antes de su distribución no solo resulta en pérdidas económicas, sino que también contribuye al aumento de los residuos en la cadena de suministro agrícola. En un contexto donde se prioriza la sostenibilidad ambiental, es crucial reducir el desperdicio y optimizar el uso de los recursos.

Asi mismo, la falta de adaptabilidad de las técnicas tradicionales frente a la creciente demanda de limones de alta calidad es otro obstáculo importante. La capacidad de los productores para competir en el mercado global se ve limitada por la ineficiencia de los métodos manuales. Un análisis de Gómez et al. (2021) sugiere que la implementación de sistemas automatizados es un paso necesario para satisfacer las expectativas de los consumidores, quienes demandan productos no solo frescos, sino también homogéneos en términos de calidad.

Frente a estos retos, la automatización, específicamente mediante el uso de visión artificial, se presenta como una solución viable. Un sistema automatizado de clasificación basado en esta tecnología tiene el potencial de superar las limitaciones del proceso manual, al aumentar la precisión y consistencia en la clasificación de los limones. La visión artificial permite que las máquinas analicen parámetros visuales como el color, tamaño y textura de los limones, garantizando una evaluación estandarizada y uniforme de la calidad del producto. Este sistema no solo reduciría el margen de error, sino que también aceleraría el proceso de clasificación, permitiendo a los productores manejar mayores volúmenes de producto en menos tiempo.

Además, la automatización reduce los costos operativos asociados a la mano de obra, lo que incrementa la rentabilidad de los productores. Al eliminar la necesidad de depender exclusivamente de trabajadores calificados, los sistemas automatizados permiten un mayor control sobre el proceso y una disminución significativa en los costos relacionados con salarios y capacitación. En este sentido, la automatización no solo mejora la eficiencia del proceso de clasificación, sino que también representa una inversión rentable para los productores a largo plazo.

\clearpage
\section{Motivación}



En el sector agrícola, la optimización de los procesos de producción es un factor clave para mejorar la eficiencia y garantizar la calidad de los productos que llegan al mercado. La clasificación de frutas, como los limones, sigue siendo un proceso que en muchos casos se realiza de forma manual, lo que no solo genera altos costos laborales, sino que también introduce variabilidad y errores humanos en la evaluación de la madurez y calidad del producto. Estas ineficiencias pueden impactar de manera negativa en la cadena de suministro, afectando tanto a los productores como a los consumidores.

El presente proyecto surge de la necesidad de automatizar este proceso utilizando tecnologías accesibles y económicas, como la ESP32-CAM, para desarrollar un sistema de clasificación basado en la visión artificial. La motivación principal es aportar una solución que no solo disminuya los tiempos y costos asociados a la selección manual de limones, sino que también mejore la precisión y consistencia de los resultados. Esto permitirá a los productores mantener estándares de calidad más altos y constantes, minimizando pérdidas y maximizando el aprovechamiento de los recursos.

Además, la implementación de un sistema de clasificación automatizado tiene un impacto directo en la reducción del desperdicio de alimentos, ya que permite clasificar los productos de manera más precisa según su estado de madurez, lo que a su vez contribuye a una mayor sostenibilidad en la producción agrícola. La utilización de la ESP32-CAM, por su accesibilidad y capacidad para realizar tareas complejas de procesamiento de imágenes, también abre la puerta para que pequeños y medianos productores puedan adoptar soluciones tecnológicas avanzadas sin necesidad de realizar grandes inversiones en infraestructura.

El prototipo pretende no solo resolver un problema específico dentro de la industria agrícola, sino también fomentar el uso de tecnologías accesibles y aplicables en diferentes áreas de la agricultura, alineándose con las tendencias actuales de automatización e innovación tecnológica.


\clearpage
\section{Justificación}

El modelo de innovación considerado para el proyecto \textit{"Automatización de Clasificación de Limones por Calidad Utilizando Visión Artificial Basada en el Color"} es la innovación incremental. 

La innovación incremental se refiere a las mejoras y ajustes realizados a productos, servicios o procesos existentes, en lugar de crear algo completamente nuevo. Esta forma de innovación busca optimizar lo que ya se tiene, aumentando la eficiencia y la efectividad de los métodos actuales.

\begin{itemize}
    \item \textbf{Mejoras en Procesos Existentes:} El proyecto se basa en la automatización de un proceso de clasificación que ya existe, el cual es tradicionalmente manual. La implementación de visión artificial y tecnología de sensores representa una mejora progresiva que busca optimizar la clasificación de limones sin revolucionar completamente el sistema agrícola.
    
    \item \textbf{Aumento de Eficiencia:} A través de esta innovación, se espera mejorar la eficiencia del proceso de clasificación, reduciendo tiempos de operación y costos, al mismo tiempo que se mantiene la esencia del proceso agrícola. La adopción de nuevas tecnologías que complementan y mejoran las prácticas existentes es un rasgo característico de la innovación incremental.
    
    \item \textbf{Reducción de Riesgos:} Al introducir mejoras graduales en lugar de cambios radicales, se minimizan los riesgos asociados con la implementación de tecnologías completamente nuevas. Esto permite a los productores adaptar sus operaciones de manera más controlada y gestionar mejor los impactos en su cadena de suministro.
    
    \item \textbf{Satisfacción del Cliente:} Las mejoras en la calidad del producto resultantes de una clasificación más precisa y eficiente pueden aumentar la satisfacción del consumidor. Esto es un objetivo clave de la innovación incremental, que busca realizar ajustes que beneficien a los usuarios finales sin alterar drásticamente la oferta existente.
    
    \item \textbf{Facilitación de la Adopción:} La innovación incremental tiende a ser más fácilmente aceptada por los usuarios, ya que se construye sobre lo que ya conocen. Esto permite una transición más suave hacia la automatización, ya que los trabajadores agrícolas y otros involucrados en la cadena de suministro pueden aprender y adaptarse a los nuevos procesos sin experimentar un cambio drástico en su entorno de trabajo.
\end{itemize}


\clearpage
\section{Alcances y limitaciones}

Para este proyecto se utilizarán diversas herramientas, tecnologías, hardware y software que permitirán una integración efectiva del sistema.

\subsection{Hardware}
\begin{itemize}
    \item \textbf{ESP32-CAM:} Este módulo se utilizará para capturar imágenes de los limones en tiempo real. Su integración se realizará mediante conexión Wi-Fi, permitiendo la transmisión de datos y la comunicación con un servidor o aplicación de procesamiento de imágenes.
    
    \item \textbf{Sensor de Color TCS3200:} Este sensor permitirá medir los parámetros de color de la piel de los limones. Se integrará directamente con el ESP32-CAM, enviando datos sobre el color detectado para su análisis.
    
    \item \textbf{Plataforma de Prueba (PCB):} Se utilizará una placa de circuito impreso personalizada para montar el ESP32-CAM y el sensor de color. Esto garantizará una conexión estable y facilitará el manejo del hardware en entornos de producción.
\end{itemize}

\subsection{Software}
\begin{itemize}
    \item \textbf{Python:} Se utilizará como lenguaje de programación principal para el procesamiento de imágenes y la implementación de algoritmos de clasificación. Python es ideal por su amplia variedad de bibliotecas y su compatibilidad con tecnologías de visión artificial.
    
    \item \textbf{OpenCV:} Esta biblioteca de visión por computadora será fundamental para procesar las imágenes capturadas por el ESP32-CAM. Se encargará de realizar la segmentación y análisis de color, así como de aplicar algoritmos de clasificación.
    
    \item \textbf{Flask:} Este framework de microservicio se utilizará para crear una API que permita la comunicación entre el hardware y el software. La API recibirá datos del sensor y de la cámara, procesará la información y devolverá los resultados de clasificación.
\end{itemize}

\subsection{Tecnologías y APIs}
\begin{itemize}
    \item \textbf{RESTful API:} Se implementará una API RESTful para permitir la comunicación entre el módulo ESP32-CAM, el servidor y cualquier aplicación frontend que muestre los resultados. Esta API será responsable de recibir imágenes y datos del sensor, procesarlos y devolver la clasificación de calidad.
    
    \item \textbf{Firebase:} Se utilizará Firebase como base de datos en tiempo real para almacenar los resultados de la clasificación y permitir el acceso a los datos desde diferentes dispositivos. Firebase facilitará la sincronización de datos entre el servidor y las aplicaciones que utilicen la información.
\end{itemize}

\subsection{Integración}
\begin{itemize}
    \item \textbf{Comunicación en Tiempo Real:} El ESP32-CAM enviará imágenes y datos del sensor de color a la API RESTful implementada en Flask. Esta API procesará la información utilizando OpenCV y devolverá la clasificación de calidad de los limones.
    
    \item \textbf{Almacenamiento y Acceso a Datos:} Los resultados de clasificación se almacenarán en Firebase, permitiendo su acceso desde aplicaciones móviles o web. Esto facilitará la consulta de datos históricos y el monitoreo de la calidad de los limones clasificados.
    
    \item \textbf{Interfaz de Usuario:} Se desarrollará una interfaz web sencilla que se conectará a la API para mostrar los resultados de clasificación en tiempo real. Los usuarios podrán visualizar los datos almacenados en Firebase, facilitando la toma de decisiones en la cadena de suministro.
\end{itemize}







%-----------------------------------------------------------------------------------------------------------------
% REFERENCIAS


\clearpage
%Let's cite! The Einstein's journal paper \cite{dirac} and the Dirac's 
%book \cite{einstein} are physics related items. 

%\Urlmuskip=0mu plus 1mu\relax
\addcontentsline{toc}{section}{Referencias} 
\printbibliography
 
\end{document}
